% Part: computability
% Chapter: recursive-functions
% Section: pr-functions-computable

\documentclass[../../../include/open-logic-section]{subfiles}

\begin{document}

\olfileid[fr]{cmp}{rec}{cmp}
\olsection{Les fonctions récursives primitives sont calculables}

Supposons qu'une fonction~$h$ soit définie par récurrence primitive :
\begin{eqnarray*}
h(\vec x, 0)     & = & f(\vec x) \\
h(\vec x, y + 1) & = & g(\vec x, y, h(\vec x, y))
\end{eqnarray*}
et que les fonctions $f$ et $g$ soient calculables. (Nous abrégeons
par~$\vec x$ la suite $x_0,\dots,x_{k-1}$.) Alors
$h(\vec x,0)$ est calculable : ce n'est autre que $f(\vec x)$,
que nous supposons calculable. On peut ensuite calculer
$h(\vec x,1)$ puisque $1=0+1$ et donc
\begin{align*}
 h(\vec x, 1) & = g(\vec x, 0, h(\vec x, 0)) =  g(\vec x, 0, f(\vec x)).
\intertext{En poursuivant de la même façon, on calcule}
h(\vec x, 2) & = g(\vec x, 1, h(\vec x, 1)) = g(\vec x, 1, g(\vec x, 0, f(\vec x)))\\
h(\vec x, 3) & = g(\vec x, 2, h(\vec x, 2)) = g(\vec x, 2, g(\vec x, 1, g(\vec x, 0, f(\vec x))))\\
h(\vec x, 4) & = g(\vec x, 3, h(\vec x, 3)) = g(\vec x, 3, g(\vec x, 2, g(\vec x, 1, g(\vec x, 0, f(\vec x)))))\\
& \vdots
\end{align*}
Pour calculer $h(\vec x,y)$ en général, il suffit donc de
calculer successivement $h(\vec x,0)$, $h(\vec x,1)$,
$\dots$, jusqu'à $h(\vec x,y)$.

Ainsi, si les fonctions $f$ et~$g$ sont calculables, une
définition par récurrence primitive fournit une nouvelle
fonction calculable. La composition de fonctions calculables
fournit elle aussi une fonction calculable, lorsque $f$ et
les fonctions $g_i$ le sont.

Les fonctions de base $\Zero$, $\Succ$ et $\Proj{n}{i}$
sont calculables. Puisque la composition et la récurrence
primitive préservent la calculabilité, toute fonction récursive
primitive est donc calculable.

\end{document}
